\begin{adjustwidth}{-\extralength}{0cm}
\reftitle{References}

\begin{thebibliography}{99}
\bibitem{AllenCahn} Allen, S.M.; Cahn, J.W. A microscopic theory for antiphase boundary motion and its application to antiphase domain coarsening. \textit{Acta Metall.} \textbf{1979}, \textit{27}, 1085--1095. https://doi.org/10.1016/0001-6160(79)90196-2.
\bibitem{Du_review} Du, Q.; Feng, X. The phase field method for geometric moving interfaces and their numerical approximations. In \textit{Handbook of Numerical Analysis}; Bonito, A., Nochetto, R.H., Eds.; Elsevier: Amsterdam, The Netherlands, 2020; Volume 21, pp.\ 425--508. https://doi.org/10.1016/bs.hna.2019.05.001.
\bibitem{Baishemirov_Symmetry} Berdyshev, A.S.; Aloev, R.D.; Bliyeva, D.N.; Dadabayev, S.; Baishemirov, Zh.D. Stability analysis of an upwind difference splitting scheme for two-dimensional Saint--Venant equations. \textit{Symmetry} \textbf{2022}, \textit{14}, 1986. https://doi.org/10.3390/sym14101986.
\bibitem{Baigereyev_Math} Baigereyev, D.R.; Alimbekova, N.B.; Berdyshev, A.S.; Madiyarov, M.N. Convergence analysis of a numerical method for a fractional model of fluid flow in fractured porous media. \textit{Mathematics} \textbf{2021}, \textit{9}, 2179. https://doi.org/10.3390/math9182179.
\bibitem{Shiyapov2026} Shiyapov, K.; Abdiramanov, Z.; Issa, Z.; Zhumaseyitova, A. High-order spectral scheme with structure maintenance and fast memory algorithm for nonlocal nonlinear diffusion equations. \textit{AppliedMath} \textbf{2026}, \textit{6}, 54. https://doi.org/10.3390/appliedmath6040054.
\bibitem{Raissi} Raissi, M.; Perdikaris, P.; Karniadakis, G.E. Physics-informed neural networks: A deep learning framework for solving forward and inverse problems involving nonlinear partial differential equations. \textit{J. Comput. Phys.} \textbf{2019}, \textit{378}, 686--707. https://doi.org/10.1016/j.jcp.2018.10.045.
\bibitem{Karniadakis2021} Karniadakis, G.E.; Kevrekidis, I.G.; Lu, L.; Perdikaris, P.; Wang, S.; Yang, L. Physics-informed machine learning. \textit{Nat. Rev. Phys.} \textbf{2021}, \textit{3}, 422--440. https://doi.org/10.1038/s42254-021-00314-5.
\bibitem{Yang_PINN} Yang, L.; Meng, X.; Karniadakis, G.E. B-PINNs: Bayesian physics-informed neural networks for forward and inverse PDE problems with noisy data. \textit{J. Comput. Phys.} \textbf{2021}, \textit{425}, 109913. https://doi.org/10.1016/j.jcp.2020.109913.
\bibitem{Lu_DeepXDE} Lu, L.; Meng, X.; Mao, Z.; Karniadakis, G.E. DeepXDE: A deep learning library for solving differential equations. \textit{SIAM Rev.} \textbf{2021}, \textit{63}, 208--228. https://doi.org/10.1137/19M1274067.
\bibitem{Khanra2026} Khanra, S.; Kukreja, V.K.; Bala, I. Physics-informed neural networks for differential equation solutions: A comprehensive review. \textit{Neurocomputing} \textbf{2026}, \textit{680}, 133317. https://doi.org/10.1016/j.neucom.2026.133317.
\bibitem{Pateras2023} Pateras, J.; Rana, P.; Ghosh, P. A taxonomic survey of physics-informed machine learning. \textit{Appl. Sci.} \textbf{2023}, \textit{13}, 6892. https://doi.org/10.3390/app13126892.
\bibitem{Lu_DeepONet} Lu, L.; Jin, P.; Pang, G.; Zhang, Z.; Karniadakis, G.E. Learning nonlinear operators via DeepONet based on the universal approximation theorem of operators. \textit{Nat. Mach. Intell.} \textbf{2021}, \textit{3}, 218--229. https://doi.org/10.1038/s42256-021-00302-5.
\bibitem{Li_DeepONetExt2025} Li, H.; Miao, Y.; Khodaei, Z.S.; Aliabadi, M.H. An architectural analysis of DeepONet and a general extension of the physics-informed DeepONet model on solving nonlinear parametric partial differential equations. \textit{Neurocomputing} \textbf{2025}, \textit{611}, 128675. https://doi.org/10.1016/j.neucom.2024.128675.
\bibitem{Li_FNO} Li, Z.; Kovachki, N.; Azizzadenesheli, K.; Liu, B.; Bhattacharya, K.; Stuart, A.; Anandkumar, A. Fourier neural operator for parametric partial differential equations. In Proceedings of the International Conference on Learning Representations (ICLR), Virtual, 3--7 May 2021. https://doi.org/10.48550/arXiv.2010.08895.
\bibitem{Wight} Wight, C.L.; Zhao, J. Solving Allen--Cahn and Cahn--Hilliard equations using the adaptive physics informed neural networks. \textit{Commun. Comput. Phys.} \textbf{2021}, \textit{29}, 930--954. https://doi.org/10.4208/cicp.OA-2020-0086.
\bibitem{Wang_NTK2022} Wang, S.; Yu, X.; Perdikaris, P. When and why PINNs fail to train: A neural tangent kernel perspective. \textit{J. Comput. Phys.} \textbf{2022}, \textit{449}, 110768. https://doi.org/10.1016/j.jcp.2021.110768.
\bibitem{Wang_causal} Wang, S.; Sankaran, S.; Perdikaris, P. Respecting causality for training physics-informed neural networks. \textit{Comput. Methods Appl. Mech. Eng.} \textbf{2024}, \textit{421}, 116813. https://doi.org/10.1016/j.cma.2024.116813.
\bibitem{Mattey} Mattey, R.; Ghosh, S. A novel sequential method to train physics-informed neural networks for Allen--Cahn and Cahn--Hilliard equations. \textit{Comput. Methods Appl. Mech. Eng.} \textbf{2022}, \textit{390}, 114474. https://doi.org/10.1016/j.cma.2021.114474.
\bibitem{Liu_FourierBias2025} Liu, Y.; Gu, H.; Yu, X.; Qin, P. Diminishing spectral bias in physics-informed neural networks using spatially-adaptive Fourier feature encoding. \textit{Neural Netw.} \textbf{2025}, \textit{182}, 106886. https://doi.org/10.1016/j.neunet.2024.106886.
\bibitem{Jagtap_XPINN2020} Jagtap, A.D.; Karniadakis, G.E. Extended physics-informed neural networks (XPINNs): A generalized space-time domain decomposition based deep learning framework for nonlinear partial differential equations. \textit{Commun. Comput. Phys.} \textbf{2020}, \textit{28}, 2002--2041. https://doi.org/10.4208/cicp.OA-2020-0164.
\bibitem{Moseley_FBPINN2023} Moseley, B.; Markham, A.; Nissen-Meyer, T. Finite basis physics-informed neural networks (FBPINNs): a scalable domain decomposition approach for solving differential equations. \textit{Adv. Comput. Math.} \textbf{2023}, \textit{49}, 62. https://doi.org/10.1007/s10444-023-10065-9.
\bibitem{Zhang_SPINN2023} Zhang, Z.-Y.; Zhang, H.; Zhang, L.-S.; Guo, L.-L. Enforcing continuous symmetries in physics-informed neural network for solving forward and inverse problems of partial differential equations. \textit{J. Comput. Phys.} \textbf{2023}, \textit{492}, 112415. https://doi.org/10.1016/j.jcp.2023.112415.
\bibitem{Zhang_SymGroup2023} Zhang, Z.-Y.; Cai, S.-J.; Zhang, H. A symmetry group based supervised learning method for solving partial differential equations. \textit{Comput. Methods Appl. Mech. Eng.} \textbf{2023}, \textit{414}, 116181. https://doi.org/10.1016/j.cma.2023.116181.
\bibitem{Brandstetter2022} Brandstetter, J.; Welling, M.; Worrall, D.E. Lie point symmetry data augmentation for neural PDE solvers. In Proceedings of the International Conference on Machine Learning (ICML); \textit{Proc. Mach. Learn. Res.} \textbf{2022}, \textit{162}, pp.\ 2241--2256. Available online: https://proceedings.mlr.press/v162/brandstetter22a/brandstetter22a.pdf (accessed on 14 September 2026).
\bibitem{Arora2024} Arora, S.; Bihlo, A.; Valiquette, F. Invariant physics-informed neural networks for ordinary differential equations. \textit{J. Mach. Learn. Res.} \textbf{2024}, \textit{25}, 233. Available online: https://jmlr.org/papers/v25/23-1511.html (accessed on 12 October 2026).
\bibitem{Liu_SymDD2024} Liu, Y.; Li, J.-Y.; Zhang, L.-S.; Guo, L.-L.; Zhang, Z.-Y. Symmetry group based domain decomposition to enhance physics-informed neural networks for solving partial differential equations. \textit{Chaos Solitons Fractals} \textbf{2024}, \textit{189}, 115658. https://doi.org/10.1016/j.chaos.2024.115658.
\bibitem{Zhu_SPINN2022} Zhu, W.; Khademi, W.; Charalampidis, E.G.; Kevrekidis, P.G. Neural networks enforcing physical symmetries in nonlinear dynamical lattices: The case example of the Ablowitz--Ladik model. \textit{Physica D} \textbf{2022}, \textit{434}, 133264. https://doi.org/10.1016/j.physd.2022.133264.
\bibitem{Shen_SAV} Shen, J.; Xu, J.; Yang, J. The scalar auxiliary variable (SAV) approach for gradient flows. \textit{J. Comput. Phys.} \textbf{2018}, \textit{353}, 407--416. https://doi.org/10.1016/j.jcp.2017.10.021.
\bibitem{Eyre} Eyre, D.J. Unconditionally gradient stable time marching the Cahn--Hilliard equation. In \textit{Computational and Mathematical Models of Microstructural Evolution}; Bullard, J.W., Chen, L.-Q., Kalia, R.K., Stoneham, A.M., Eds.; MRS Symposium Proceedings; Materials Research Society: Warrendale, PA, USA, 1998; Volume 529, pp.\ 39--46. https://doi.org/10.1557/PROC-529-39.
\bibitem{Huang_GSAV} Huang, F.; Shen, J.; Yang, Z. A highly efficient and accurate new scalar auxiliary variable approach for gradient flows. \textit{SIAM J. Sci. Comput.} \textbf{2020}, \textit{42}, A2514--A2536. https://doi.org/10.1137/19M1298627.
\bibitem{Jiang_RSAV} Jiang, M.; Zhang, Z.; Zhao, J. Improving the accuracy and consistency of the scalar auxiliary variable (SAV) method with relaxation. \textit{J. Comput. Phys.} \textbf{2022}, \textit{456}, 110954. https://doi.org/10.1016/j.jcp.2022.110954.
\bibitem{Zhang_ETD2024} Zhang, H.; Liu, L.; Qian, X.; Song, S. Quantifying and eliminating the time delay in stabilization exponential time differencing Runge--Kutta schemes for the Allen--Cahn equation. \textit{ESAIM Math. Model. Numer. Anal.} \textbf{2024}, \textit{58}, 191--221. https://doi.org/10.1051/m2an/2023101.
\bibitem{Gao_pIFRK2023} Gao, Z.; Zhang, H.; Qian, X.; Song, S. High-order unconditionally maximum-principle-preserving parametric integrating factor Runge--Kutta schemes for the nonlocal Allen--Cahn equation. \textit{Appl. Numer. Math.} \textbf{2023}, \textit{194}, 97--114. https://doi.org/10.1016/j.apnum.2023.08.010.
\bibitem{Du_MPPDG2022} Du, J.; Chung, E.; Yang, Y. Maximum-principle-preserving local discontinuous Galerkin methods for Allen--Cahn equations. \textit{Commun. Appl. Math. Comput.} \textbf{2022}, \textit{4}, 353--379. https://doi.org/10.1007/s42967-020-00118-x.
\bibitem{Feng_MPPFD2021} Feng, J.; Zhou, Y.; Hou, T. A maximum-principle preserving and unconditionally energy-stable linear second-order finite difference scheme for Allen--Cahn equations. \textit{Appl. Math. Lett.} \textbf{2021}, \textit{118}, 107179. https://doi.org/10.1016/j.aml.2021.107179.
\bibitem{Kim_MPP3D2025} Kim, J.; Hwang, Y. Maximum principle-preserving computational algorithm for the 3D high-order Allen--Cahn equation. \textit{Mathematics} \textbf{2025}, \textit{13}, 1085. https://doi.org/10.3390/math13071085.
\bibitem{ICCMA2026abstract} Nauryzbay, A.; Abdiramanov, Z.; Zhumazhanov, E.; Urnek, D.; Ertai, K. Investigation of Physics-Informed Neural Networks for Solving the Allen--Cahn Equation. In Proceedings of the 4th International Conference: Constructive Mathematical Analysis (ICCMA 2026), Konya, T\"urkiye, 1--2 July 2026; Abstract Book, p.~4. Available online: https://iccma2026.selcuk.edu.tr/wp-content/uploads/2026/07/BookofAbstract.pdf (accessed on 14 September 2026).
\bibitem{Trefethen} Trefethen, L.N. \textit{Spectral Methods in MATLAB}; SIAM: Philadelphia, PA, USA, 2000. https://doi.org/10.1137/1.9780898719598.
\bibitem{ShenTangWang} Shen, J.; Tang, T.; Wang, L.-L. \textit{Spectral Methods: Algorithms, Analysis and Applications}; Springer: Berlin/Heidelberg, Germany, 2011. https://doi.org/10.1007/978-3-540-71041-7.
\bibitem{Mostajeran_cPIKAN2025} Mostajeran, F.; Faroughi, S.A. Scaled-cPIKANs: Spatial variable and residual scaling in Chebyshev-based physics-informed Kolmogorov--Arnold networks. \textit{J. Comput. Phys.} \textbf{2025}, \textit{537}, 114116. https://doi.org/10.1016/j.jcp.2025.114116.
\bibitem{Guo_EnergySeq2024} Guo, J.; Wang, H.; Hou, C. An adaptive energy-based sequential method for training PINNs to solve gradient flow equations. \textit{Appl. Math. Comput.} \textbf{2024}, \textit{479}, 128890. https://doi.org/10.1016/j.amc.2024.128890.
\bibitem{DeMouraUllrich2021} de Moura, L.; Ullrich, S. The Lean 4 theorem prover and programming language. In \textit{Automated Deduction -- CADE 28}; Platzer, A., Sutcliffe, G., Eds.; Lecture Notes in Computer Science; Springer: Cham, Switzerland, 2021; Volume 12699, pp.\ 625--635. https://doi.org/10.1007/978-3-030-79876-5\_37.
\bibitem{Bobbin2023} Bobbin, M.P.; Sharlin, S.; Feyzishendi, P.; Dang, A.H.; Wraback, C.M.; Josephson, T.R. Formalizing chemical physics using the Lean theorem prover. \textit{Digit. Discov.} \textbf{2024}, \textit{3}, 264--280. https://doi.org/10.1039/D3DD00077J.
\bibitem{Ugwuanyi2025} Ugwuanyi, E.D.; Jones, C.T.; Velkey, J.; Josephson, T.R. Benchmarking energy calculations using formal proofs. \textit{Mol. Phys.} \textbf{2025}, \textit{123}, e2539421. https://doi.org/10.1080/00268976.2025.2539421.
\bibitem{Tsachouridis2020} Tsachouridis, V.A.; Giantamidis, G.; Basagiannis, S.; Kouramas, K. Formal analysis of the Schulz matrix inversion algorithm: A paradigm towards computer aided verification of general matrix flow solvers. \textit{Numer. Algebra Control Optim.} \textbf{2020}, \textit{10}, 177--206. https://doi.org/10.3934/naco.2019047.
\bibitem{ref_Xia2023} Xia, M.; B\"ottcher, L.; Chou, T. Spectrally adapted physics-informed neural networks for solving unbounded domain problems. \textit{Mach. Learn.: Sci. Technol.} \textbf{2023}, \textit{4}, 025024. https://doi.org/10.1088/2632-2153/acd0a1.
\bibitem{ref_Kutuk2025} K\"ut\"uk, M.; Y\"ucel, H. Energy dissipation preserving physics informed neural network for Allen--Cahn equations. \textit{J. Comput. Sci.} \textbf{2025}, \textit{87}, 102577. https://doi.org/10.1016/j.jocs.2025.102577.
\bibitem{ref_He2019} He, D.; Pan, K. Maximum norm error analysis of an unconditionally stable semi-implicit scheme for multi-dimensional Allen--Cahn equations. \textit{Numer. Methods Partial Differ. Equ.} \textbf{2019}, \textit{35}, 955--975. https://doi.org/10.1002/num.22333.
\bibitem{ref_Du2024} Du, Y.; Chalapathi, N.; Krishnapriyan, A. Neural spectral methods: Self-supervised learning in the spectral domain. In Proceedings of the 12th International Conference on Learning Representations (ICLR), Vienna, Austria, 7--11 May 2024. https://doi.org/10.48550/arXiv.2312.05225.
\bibitem{ref_Bizzi2025} Bizzi, A.; Moreira, L.M.; Marques, M.; Mendon\c{c}a, L.; de Oliveira, C.J.; Balestro, V.; dos Santos Fernandez, L.; Yukimura, D.; Petrov, P.; Pereira, J.M.; Novello, T.; Nissenbaum, L. Neuro-spectral architectures for causal physics-informed networks. In \textit{Advances in Neural Information Processing Systems 38 (NeurIPS 2025)}; 2025; pp.\ 60600--60630. https://doi.org/10.52202/085713-1818.
\bibitem{ref_Zhang2024} Zhang, J.; Zhang, S.; Shen, J.; Lin, G. Energy-dissipative evolutionary deep operator neural networks. \textit{J. Comput. Phys.} \textbf{2024}, \textit{498}, 112638. https://doi.org/10.1016/j.jcp.2023.112638.
\bibitem{ref_Lin2025} Lin, G.; Mou, C.; Zhang, J. Energy-dissipative evolutionary Kolmogorov--Arnold networks for complex PDE systems. \textit{J. Comput. Phys.} \textbf{2025}, \textit{541}, 114326. https://doi.org/10.1016/j.jcp.2025.114326.
\bibitem{ref_Du2021b} Du, Q.; Ju, L.; Li, X.; Qiao, Z. Maximum bound principles for a class of semilinear parabolic equations and exponential time-differencing schemes. \textit{SIAM Rev.} \textbf{2021}, \textit{63}, 317--359. https://doi.org/10.1137/19M1243750.
\bibitem{ref_Tang2016} Tang, T.; Yang, J. Implicit-explicit scheme for the Allen--Cahn equation preserves the maximum principle. \textit{J. Comput. Math.} \textbf{2016}, \textit{34}, 451--461. https://doi.org/10.4208/jcm.1603-m2014-0017.
\bibitem{ref_Grossmann2024} Grossmann, T.G.; Komorowska, U.J.; Latz, J.; Sch\"onlieb, C.-B. Can physics-informed neural networks beat the finite element method? \textit{IMA J. Appl. Math.} \textbf{2024}, \textit{89}, 143--174. https://doi.org/10.1093/imamat/hxae011.
\bibitem{ref_Tekriwal2021} Tekriwal, M.; Duraisamy, K.; Jeannin, J.-B. A formal proof of the Lax equivalence theorem for finite difference schemes. In \textit{NASA Formal Methods}; Dutle, A., Moscato, M.M., Titolo, L., Mu\~noz, C.A., Perez, I., Eds.; Lecture Notes in Computer Science; Springer: Cham, Switzerland, 2021; Volume 12673, pp.\ 322--339. https://doi.org/10.1007/978-3-030-76384-8\_20.
\bibitem{ref_Zhao2021} Zhao, J. Discovering phase field models from image data with the pseudo-spectral physics informed neural networks. \textit{Commun. Appl. Math. Comput.} \textbf{2021}, \textit{3}, 357--369. https://doi.org/10.1007/s42967-020-00105-2.
\bibitem{ref_Bruna2024} Bruna, J.; Peherstorfer, B.; Vanden-Eijnden, E. Neural Galerkin schemes with active learning for high-dimensional evolution equations. \textit{J. Comput. Phys.} \textbf{2024}, \textit{496}, 112588. https://doi.org/10.1016/j.jcp.2023.112588.
\bibitem{ref_Du2021a} Du, Y.; Zaki, T.A. Evolutional deep neural network. \textit{Phys. Rev. E} \textbf{2021}, \textit{104}, 045303. https://doi.org/10.1103/PhysRevE.104.045303.
\bibitem{ref_Cooley2025} Cooley, M.; Shankar, V.; Kirby, R.M.; Zhe, S. Fourier PINNs: From strong boundary conditions to adaptive Fourier bases. \textit{Trans. Mach. Learn. Res.} \textbf{2025}. https://doi.org/10.48550/arXiv.2410.03496.
\bibitem{ref_Li2026} Li, C.; Yu, W.; Wang, Q. Energy dissipation rate guided adaptive sampling for physics-informed neural networks: Resolving surface-bulk dynamics in Allen--Cahn systems. \textit{Int. J. Numer. Anal. Model.} \textbf{2026}, \textit{23}, 576--604. https://doi.org/10.4208/ijnam2026-1025.
\bibitem{ref_Geng2025} Geng, Y.; Burkovska, O.; Ju, L.; Zhang, G.; Gunzburger, M. An end-to-end deep learning method for solving nonlocal Allen--Cahn and Cahn--Hilliard phase-field models. \textit{Comput. Methods Appl. Mech. Eng.} \textbf{2025}, \textit{436}, 117721. https://doi.org/10.1016/j.cma.2024.117721.
\bibitem{ref_Shen2016} Shen, J.; Tang, T.; Yang, J. On the maximum principle preserving schemes for the generalized Allen--Cahn equation. \textit{Commun. Math. Sci.} \textbf{2016}, \textit{14}, 1517--1534. https://doi.org/10.4310/CMS.2016.v14.n6.a3.
\bibitem{ref_Dong2021} Dong, S.; Ni, N. A method for representing periodic functions and enforcing exactly periodic boundary conditions with deep neural networks. \textit{J. Comput. Phys.} \textbf{2021}, \textit{435}, 110242. https://doi.org/10.1016/j.jcp.2021.110242.
\bibitem{ref_Sukumar2022} Sukumar, N.; Srivastava, A. Exact imposition of boundary conditions with distance functions in physics-informed deep neural networks. \textit{Comput. Methods Appl. Mech. Eng.} \textbf{2022}, \textit{389}, 114333. https://doi.org/10.1016/j.cma.2021.114333.
\bibitem{ref_Lu2021} Lu, L.; Pestourie, R.; Yao, W.; Wang, Z.; Verdugo, F.; Johnson, S.G. Physics-informed neural networks with hard constraints for inverse design. \textit{SIAM J. Sci. Comput.} \textbf{2021}, \textit{43}, B1105--B1132. https://doi.org/10.1137/21M1397908.
\bibitem{ref_CardosoBihlo2025} Cardoso-Bihlo, E.; Bihlo, A. Exactly conservative physics-informed neural networks and deep operator networks for dynamical systems. \textit{Neural Netw.} \textbf{2025}, \textit{181}, 106826. https://doi.org/10.1016/j.neunet.2024.106826.
\bibitem{ref_Krishnapriyan2021} Krishnapriyan, A.S.; Gholami, A.; Zhe, S.; Kirby, R.M.; Mahoney, M.W. Characterizing possible failure modes in physics-informed neural networks. In \textit{Advances in Neural Information Processing Systems 34 (NeurIPS 2021)}; 2021; pp.\ 26548--26560. https://doi.org/10.48550/arXiv.2109.01050.
\bibitem{ref_Rathore2024} Rathore, P.; Lei, W.; Frangella, Z.; Lu, L.; Udell, M. Challenges in training PINNs: A loss landscape perspective. In Proceedings of the 41st International Conference on Machine Learning (ICML); \textit{Proc. Mach. Learn. Res.} \textbf{2024}, \textit{235}, pp.\ 42159--42191. https://doi.org/10.48550/arXiv.2402.01868.
\bibitem{ref_Wang2021a} Wang, S.; Teng, Y.; Perdikaris, P. Understanding and mitigating gradient flow pathologies in physics-informed neural networks. \textit{SIAM J. Sci. Comput.} \textbf{2021}, \textit{43}, A3055--A3081. https://doi.org/10.1137/20M1318043.
\bibitem{ref_Rahaman2019} Rahaman, N.; Baratin, A.; Arpit, D.; Draxler, F.; Lin, M.; Hamprecht, F.A.; Bengio, Y.; Courville, A. On the spectral bias of neural networks. In Proceedings of the 36th International Conference on Machine Learning (ICML); \textit{Proc. Mach. Learn. Res.} \textbf{2019}, \textit{97}, pp.\ 5301--5310. https://doi.org/10.48550/arXiv.1806.08734.
\bibitem{ref_McGreivy2024} McGreivy, N.; Hakim, A. Weak baselines and reporting biases lead to overoptimism in machine learning for fluid-related partial differential equations. \textit{Nat. Mach. Intell.} \textbf{2024}, \textit{6}, 1256--1269. https://doi.org/10.1038/s42256-024-00897-5.
\bibitem{ref_Mishra2023} Mishra, S.; Molinaro, R. Estimates on the generalization error of physics-informed neural networks for approximating PDEs. \textit{IMA J. Numer. Anal.} \textbf{2023}, \textit{43}, 1--43. https://doi.org/10.1093/imanum/drab093.
\bibitem{ref_Boldo2013} Boldo, S.; Cl\'ement, F.; Filli\^atre, J.-C.; Mayero, M.; Melquiond, G.; Weis, P. Wave equation numerical resolution: A comprehensive mechanized proof of a C program. \textit{J. Autom. Reason.} \textbf{2013}, \textit{50}, 423--456. https://doi.org/10.1007/s10817-012-9255-4.
\bibitem{ref_Boldo2017} Boldo, S.; Cl\'ement, F.; Faissole, F.; Martin, V.; Mayero, M. A Coq formal proof of the Lax--Milgram theorem. In Proceedings of the 6th ACM SIGPLAN Conference on Certified Programs and Proofs (CPP); ACM: New York, NY, USA, 2017; pp.\ 79--89. https://doi.org/10.1145/3018610.3018625.
\bibitem{ref_Immler2018} Immler, F. A verified ODE solver and the Lorenz attractor. \textit{J. Autom. Reason.} \textbf{2018}, \textit{61}, 73--111. https://doi.org/10.1007/s10817-017-9448-y.
\end{thebibliography}
\end{adjustwidth}
\end{document}

